% Generated draft: requires mathematical and source review.
\documentclass[11pt]{article}
\usepackage[margin=25mm]{geometry}
\usepackage{fontspec}
\setmainfont{Latin Modern Roman}
\usepackage{amsmath,amssymb,mathtools,unicode-math}
\setmathfont{Latin Modern Math}
\usepackage{enumitem,needspace,xurl,hyperref}
\setlength{\parindent}{0pt}
\setlength{\parskip}{5pt}
\setlength{\emergencystretch}{3em}
\newcommand{\N}{\mathbb{N}}
\newcommand{\Z}{\mathbb{Z}}
\newcommand{\Q}{\mathbb{Q}}
\newcommand{\R}{\mathbb{R}}
\newcommand{\C}{\mathbb{C}}
\newcommand{\F}{\mathbb{F}}
\newcommand{\A}{\mathbb{A}}
\newcommand{\PP}{\mathbb{P}}
\newcommand{\E}{\mathbb{E}}
\newcommand{\eps}{\varepsilon}
\newcommand{\dd}{\,\mathrm{d}}
\begin{document}
\section*{UnsolvedMath ID: 11000064}
\textbf{Source problem code:} catalogue entry
\subsection{Definitions and mathematical statement}
Let $F_n$ be the free group of rank $n$, and let
\[
 \operatorname{IA}_n=\ker\bigl(\operatorname{Aut}(F_n)
 \longrightarrow\operatorname{GL}(H_1(F_n;\Z))\bigr),
 \qquad H_1(F_n;\Z)\cong\Z^n.
\]
Problem 5.20 asks for an explicit finite presentation of the Malcev Lie algebra $L(\operatorname{IA}_n)$, as an analogue of the result attributed to Hain for a Torelli group.

In the standard rational Malcev interpretation this is the pronilpotent Lie algebra of the rational unipotent completion. A finite presentation would give finitely many generators and relators in a completed free Lie algebra over $\Q$, with the relators generating a closed ideal. The excerpt does not state whether $L$ instead denotes an associated graded Lie algebra, nor its completion and grading conventions; a finite presentation in the ordinary algebraic category is a different requirement. The intended rank range and any uniformity in $n$ are also unstated.

The extracted following discussion concerns information at the prime $2$ for $\mathcal I_{g,1}$. It says that the mod-$2$ Lie algebra $L_2(\mathcal I_{g,1})$ is substantially larger than the part attributed to $L_0(\mathcal I_{g,1})$, and relates a $2$-torsion object $B_2$ to the Rochlin invariant rather than the purely algebraic Johnson-homomorphism contribution. It contrasts this with $\operatorname{IA}_n$, whose corresponding Johnson homomorphism is reported to account for the first-degree abelian information. The definitions of $L_0,L_2,B_2$ and referenced equation (16) are absent. In particular, rational Lie algebras have no direct reduction modulo $2$ without an integral form, and integral $H^1(G;\Z)$ is torsion-free; the excerpt's conflicting homology/cohomology notation cannot justify a claim of integral first-cohomology torsion.
\subsection{Short English statement}
Give explicit finite generators and relations for the rational Malcev Lie algebra of the free-group Torelli subgroup $\operatorname{IA}_n$, after fixing its completion convention. The adjoining mod-$2$ discussion has unresolved integral-form and homology-notation issues.
\subsection{Sources}
\begin{itemize}
\item Local registry: \url{data/unsolvedmath/problems.json}, permanent UnsolvedMath ID $11000064$, exact source code \texttt{catalogue entry}.
\item Farb, \emph{Problems on Mapping Class Groups and Related Topics}, Problem 5.20, PDF page 51: \url{https://www.math.uchicago.edu/~farb/papers/mcgbook.pdf}. Supplied access date: 2026-07-29; extraction: pdf-font-structure; page-disambiguated numbering among 329 headings.
\item The source invokes Hain's theorem, the Johnson homomorphism, the Rochlin invariant, and equation (16), but supplies no precise theorem statement, equation, or bibliography entry in this record. The registry describes analogous Malcev questions as active without exact references.
\end{itemize}
\end{document}
